\section{Operator construction\label{sec:theory}}

The energy of an eigenstate of the hamiltonian of a quantum field
theory can be determined by computing the correlation function between
creation and annihilation operators $\chi$ at Euclidean times $t$ and $t'$;
\begin{equation}
   C(t',t)  = \Big\langle \chi(t')\,  \chi^\dagger(t) \Big\rangle.
\end{equation}
Inserting a complete set of eigenstates of the hamiltonian, such that
$\hat{H} | k \rangle = E_k  | k\rangle$, this correlation function
decomposes into a sum of contributions from all states in the spectrum
with the same quantum numbers as the source operators,
\begin{equation}
   C(t',t)  = \sum_k \big|\langle \chi | k \rangle\big|^2 \, e^{-{E_k} (t'-t)}.
\end{equation}
In order to
measure energies of low-lying states, it it crucial to construct operators that
overlap predominantly with these light modes.
This exposes the asymptotic behaviour of the correlator at earlier
time-separations and enables more statistically accurate determinations of
energies.
There is a good deal of freedom in
the construction of appropriate operators. Once the constraints of symmetries
and temporal localisation are imposed, any function of the fields in the path
integral can be used in principle.

Smearing is a well-established means of defining an initial step in the
construction of creation operators. Rather than applying a creation operator to
the fields directly, a smoothing function is applied first. The function should
preserve as many symmetries as possible while effectively removing the
presence of short-range modes, which make an insignificant contribution to the
low-energy correlation function. In contemporary simulations, a popular form
of this operation in theories where fermion fields couple to gauge bosons is
the Jacobi smearing method \cite{Allton:1993wc}.

Gauge-covariant quark smearings based on the lattice Laplacian start
with the simplest representation of the second-order three-dimensional
differential operator
\begin{equation}
   -\nabla^2_{xy}(t) = 6 \delta_{xy} - \sum_{j=1}^3
      \left(
	  \tilde{U}_j(x,t)         \delta_{x+\hat\jmath,y}
	 +\tilde{U}^\dagger_j(x-\hat\jmath,t) \delta_{x-\hat\jmath,y}
	  \right),
\end{equation}
where the gauge fields, $\tilde{U}$ may be constructed from an appropriate
covariant gauge-field-smearing algorithm \cite{Morningstar:2003gk}.
After defining the Laplace operator, a simple smearing can be written
\begin{equation}
  J_{\sigma, n_\sigma}(t) = \left( 1 + \frac{\sigma \nabla^2(t)}{n_\sigma} \right)^{n_\sigma},
  \label{eqn:smearing}
\end{equation}
where $\sigma$ and $n_\sigma$ are tunable parameters that can be used to
optimise projection onto the states under investigation.
For large $n_\sigma$, this approximates the exponential of $\sigma\nabla^2$,
{\it i.e.}
\begin{equation}
  \lim_{n_\sigma\rightarrow\infty} J_{\sigma, n_\sigma}(t) = \exp\left(\sigma
      \nabla^2(t)\right).
\label{eqn:jacobi}
\end{equation}
The resulting exponential suppression of higher eigenmodes of the lattice
Laplace operator means only a small number of the lowest modes contribute
substantially to $J$.

This observation suggests the smearing operator can be approximated
by forming an eigenvector representation, truncated to the lowest
modes.  Since the smearing operators used in lattice calculations are
approximated well by low-rank constructions, a definition of smearing can be
chosen to enforce this more absolutely.  Let $V_M$ be the vector space
of scalar fields charged under the fundamental representation of the gauge
group on a particular time-slice. $V_M$ has rank $M = N_c \times N_x \times
N_y \times N_z$ where $N_c$ is the number of colours, and $N_x,N_y$ and $N_z$
are the extents of the lattice in the three spatial directions.
Now define smearing to be a well-chosen operator of rank $N \ll M$.  This class
of operators will be called ``distillation'' operators.  There is a substantial
benefit to doing this: if the rank of the operator is sufficiently small, all
elements of the propagation matrix from this space can be constructed at an
affordable computational cost (which is proportional to the rank of the space, $N$).
Consequently, correlation functions involving arbitrarily intricate hadron
creation operators can be measured with a fixed inversion overhead.

Define the distillation operator on time-slice $t$ as a product of an $M \times
N$ matrix and its hermitian conjugate:
\begin{equation}
 \Box(t) = V(t) V^\dagger(t)
\implies \Box_{xy}(t) = \sum_{k=1}^N v_x^{(k)} (t) v_y^{(k)\dag} (t).
    \label{eqn:box}
\end{equation}
The $k^{\rm th}$ column of $V(t)$ contains the $k^{\rm th}$ eigenvector
of $\nabla^2$ evaluated on the background of the spatial gauge fields of
time-slice $t$, once the eigenvectors have been sorted by eigenvalue. This is
the projection operator into $V_N$, the subspace spanned by these eigenmodes,
so $\Box^2 = \Box$.
When the number of eigenvectors included is the same as the dimension of
$V_M$, {\em i.e.} $N = M$, the distillation operator becomes the identity, and
fields acted upon are unsmeared.

The Laplace operator inherits many symmetries of the vacuum. It transforms
like a scalar under rotations, is covariant under gauge transformation and
is parity and charge conjugation invariant. If the action of one of these
symmetries on the Laplace operator maps $\nabla^2 $ onto $\tilde\nabla^2$, then
there is a unitary transformation, $R$ on $V_M$ such that
\beq
R \tilde\nabla^2 R^\dagger = \nabla^2.
\eeq
This implies that if $v$ is an eigenvector of $\nabla^2$ then the eigenvectors
of $\tilde\nabla^2$ are $R v$.
Considering the definition of the distillation operator given in
Eq.~\ref{eqn:box}, the transformed operator must then obey
\beq
R \stackrel{\sim}{\Box} R^\dagger = \Box,
\eeq
and so correlation functions constructed using distilled fields have the same
symmetry properties on the lattice as those constructed using Laplacian smearing
methods.

\subsection{Meson two-point correlation functions}\label{subsec:theory-mesons}

Consider the momentum-projected creation and annihilation
operators of an isovector meson, $\bar{u} \Gamma^A d$ and
$\bar{d} \Gamma^B u$, where $\Gamma$ acts in spin and color as well as
coordinate space. Applying the distillation operator $\Box$ onto each
quark field, the creation operator at three-momentum $\vec{p}$ is written as
\begin{equation}
\chi^\dagger_M(\vec{p},t)
  = \bar{u}_x(t) \Box_{xy}(t) \cdot e^{-i p\cdot y}\, \Gamma^A_{yz}(t) \cdot \Box_{zw}(t) d_w(t),
\label{eqn:meson_quark_op}
\end{equation}
where there is an implied summation over repeated spatial indices. In a shorthand notation the correlation function can be written as
\begin{equation}
  C^{(2)}_M(t',t) = \Big\langle \bar{d}(t')\Box(t')\Gamma^B(t')\Box(t')u(t')\,\cdot\,
                      \bar{u}(t)\Box(t)\Gamma^A(t)\Box(t) d(t)
    \Big\rangle.
\end{equation}
After evaluating the quark field path-integral and inserting the outer-product
definition of the distillation operator $\Box$ from Eq.~\ref{eqn:box}, the
correlator can be written
\begin{equation}
  C^{(2)}_M(t',t) = \mbox{Tr} \Bigl[
      \Phi^B(t') \,
      \tau(t',t) \,\Phi^A(t) \, \tau(t,t') \Bigr],
\end{equation}
where
\begin{equation}
  \Phi^A_{\alpha\beta}(t) = V^\dagger(t) \left[ \Gamma^A(t) \right]_{\alpha\beta} V(t)
  \equiv V^\dagger(t) {\cal D}^A(t) V(t) S^A_{\alpha\beta} ,
\label{eqn:meson_op}
\end{equation}
and
\begin{equation}
  \tau_{\alpha\beta }(t',t) = V^\dagger(t') M^{-1}_{\alpha\beta}(t',t) V(t),
\label{eqn:peram}
\end{equation}
with $M$ the lattice representation of the Dirac operator
and where the quark spin indices, $\alpha,\beta$ of $\Phi$ and $\tau$ have been
explicitly written. $\Phi$ has a well-defined momentum, while there is no
explicit momentum projection in the definition of $\tau$. Often, $\Phi$ can be
decomposed into terms that act only within coordinate and color space ${\cal
D}^A$ and only within spin space $S^{A}$. Note that $\Phi$ and $\tau$ are
square matrices of dimension $N \times N_\sigma$ where $N_\sigma$ is the number
of components in a lattice Dirac spinor. Therefore it requires just
$N \times N_\sigma$ operations of the inverse of the fermion matrix on a
vector in order to compute all elements of $\tau$, the ``perambulator''.
Notice also that the choice of source and sink operators is entirely
independent of the computation of $\tau$; any source and sink operators can be
correlated {\it a posteriori} once all elements of the $\tau$ matrix have been
computed and stored. The method straightforwardly extends to the determination
of correlation functions for mesons composed of different, non-degenerate
quark flavors.

In the determination of an isoscalar meson correlation function, evaluation of
disconnected terms is required.
The disconnected diagram can be similarly represented once distilled fields are
used in creation and annihilation operators. Such a term would
comprise two separate traces over the distillation space:
\begin{equation}
  C^{(2,\mathrm{disc})}_M(t',t) =
   \mbox{Tr}\Bigl[ \Phi^A(t) \tau(t,t) \Bigr]
   \mbox{Tr}\Bigl[ \Phi^B(t') \tau^\dagger(t' ,t') \Bigr].
\end{equation}
An exact determination of this expression requires computing $\tau(t,t)$ for
all time-slices $t$.

Since distilled single-particle operators can be projected onto definite
momentum at both source and sink, multi-meson correlators can also be
constructed using creation and annihilation operators of the form
\begin{equation}
\chi_{MM}(\vec{q}=\vec{p}_1 - \vec{p}_2; t) = \chi_M(\vec{p}_1,t)
  \chi_M(-\vec{p}_2,t),
\label{eqn:multi_meson_quark_op}
\end{equation}
where $p_1$ and $p_2$ are the three-momenta of the single particle operators
in Eq.~\ref{eqn:meson_quark_op}. After integration over the quark fields, the
resulting diagrams can again be computed by taking traces over products of
construction operators in the distillation space with the perambulators.
The precise structure of these functions depends on the quark flavor content of
the source and sink operators, so details are not presented here.  As before,
these matrices are small compared to the dimension of the space of quark
fields, and once the quark propagation is encoded in the perambulator,
these multi-hadron correlation function can be computed.

\subsection{Baryon two-point correlation functions}\label{subsec:theory-baryons}

The factorization technique can also be applied to baryons.
To illustrate the concepts involved, consider just the isospin-1/2 sector
although the technique generalizes naturally to other baryon multiplets.
An annihilation operator involving displacements
as well as coefficients in spin space is written:
\begin{equation}
\chi_B(t) = \epsilon^{abc} S_{\alpha_1\alpha_2\alpha_3}
({\cal D}_1\Box d)^a_{\alpha_1}
({\cal D}_2\Box u)^b_{\alpha_2}
({\cal D}_3\Box u)^c_{\alpha_3}(t),
\end{equation}
where the color indices of the quark fields acted upon
by the displacement operators ${\cal D}_i$ are contracted with the
antisymmetric tensor, and repeated spin indices are summed.
After integration over quark fields, the correlation function
\begin{equation}
   C^{(2)}_B(t',t) = -\Big\langle \chi_B(t') \, \bar{\chi}_B(t)\Big\rangle
\end{equation}
can be factored into perambulator terms, Eq.~\ref{eqn:peram}
as well as creation and annihilation operators, where the latter can be
written as
\begin{equation}
 \Phi^{(i,j,k)}_{\alpha_1\alpha_2\alpha_3}(t) = \epsilon^{abc}
\left({\cal D}_1 v^{(i)}\right)^{a}
\left({\cal D}_2 v^{(j)}\right)^{b}
\left({\cal D}_3 v^{(k)}\right)^{c}(t)\;
S_{\alpha_1\alpha_2\alpha_3}\ .
\label{eqn:baryon_op}
\end{equation}
The spin terms factor from the antisymmetric contraction of the vectors
$v$.

There are two terms in the resulting correlation function
that involve tensor contractions of the creation and annihilation operators
with the perambulators
\begin{align}
C^{(2)}_B[\tau_d,\tau_u,\tau_u](t',t) &=
  \Phi^{(i,j,k)}(t')
  \tau_d^{(i,\bar{i})}(t',t)
  \tau_u^{(j,\bar{j})}(t',t)
  \tau_u^{(k,\bar{k})}(t',t)
  \Phi^{(\bar{i},\bar{j},\bar{k})*}(t)
  \nonumber\\
 &\quad- \Phi^{(i,j,k)}(t')
  \tau_d^{(i,\bar{i})}(t',t)
  \tau_u^{(j,\bar{k})}(t',t)
  \tau_u^{(k,\bar{j})}(t',t)
  \Phi^{(\bar{i},\bar{j},\bar{k})*}(t),
\end{align}
where there is an implicit tensor contraction over the internal spin
indices, and where the quark labels $d$ and $u$ are used to denote
the corresponding flavors of the quark perambulator terms in
Eq.~\ref{eqn:peram}.

Similar to the meson case, the choice of source and sink operators is
independent of the computation of the perambulators $\tau_d$ and $\tau_u$.
The computation of these matrices can also be shared with computation of meson
correlators.  A baryon correlation function can be evaluated {\em a posteriori}
using the contractions of the vectors with displacements in
Eq.~\ref{eqn:baryon_op} which can also be shared among the source and sink
operators. These contractions do not involve spin components, thus making the
storage of $\Phi$ manageable.

\subsection{Meson three-point correlation functions}\label{subsec:theory-mesons-3pt}

A generic meson three-point function is written
\beq
   C^{(3)}(t_f, t, t_i) = \Big\langle
\big( \bar{d}\Box \Gamma^B \Box u\big)(t_f) \cdot
\big(\bar{u}\Gamma u\big)(t) \cdot
\big(\bar{u} \Box \Gamma^A \Box d \big)(t_i)
\Big\rangle,
\eeq
where there is no smearing in the operator on timeslice $t$.
The completely connected Wick contraction can be expressed as
\beq
 C^{(3)}_{\mathrm{conn}}(t_f, t, t_i) = \mbox{Tr}\Bigl[\Phi^B(t_f)\, J(t_f,t,t_i) \,\Phi^A(t_i)\, \gamma_5 \tau^\dag(t_f,t_i) \gamma_5\Bigr]
\eeq
where the ``generalized perambulator'' is defined by
\begin{equation}
J_{\alpha\beta}(t_f,t,t_i) =  V^\dag(t_f) \left( M^{-1}(t_f,t)\, \Gamma(t) \,
M^{-1}(t,t_i)  \right)_{\alpha\beta} V(t_i). \label{eqn:genperam}
\end{equation}
In general, a disconnected term may also appear; this can be factorised as the
product of a two-point function (evaluated according to the algorithm of
Sec.~\ref{subsec:theory-mesons}) and a disconnected insertion. Since the
disconnected trace does not involve distilled fields, the evaluation of this
insertion is not discussed further at this point.
To compute the connected three-point correlation function, there are two
distinct sets of inversions needed. The first is from the distillation vectors
at the source time-slice $t_i$ and the second from the sink time-slice $t_f$.
The generalised perambulator of Eq.~\ref{eqn:genperam} is then constructed by
contracting the two solutions from these two sets of inversions.
The current insertion is encoded in the choice of operator $\Gamma$.
Any momentum insertion at the current operator involves a Fourier transform on
each time-slice $t$.

\subsection{Baryon three-point correlation functions}\label{subsec:theory-baryons-3pt}

Baryon three-point correlation functions can also be expressed in terms of
the generalized perambulators of Eq.~\ref{eqn:genperam}.
Consider a three-point correlation function where a baryon is created on
time-slice $t_i$, then is acted on by a current operator on time-slice $t$ and
is subsequently annihilated at $t_f$:
\begin{equation}
   C^{(3)}_B(t_f, t, t_i) = - \Big\langle B(t_f)\cdot \big(\bar{u}\Gamma u \big)(t) \cdot \bar{B}(t_i) \Big\rangle.
\end{equation}
After quark-field integration, this can be written as a sum of both
a connected three-point contribution
and the product of a disconnected insertion and a two-point correlator.
As before, the discussion here is restricted to isospin-1/2 light baryons
although the technique is quite general.
Note that the disconnected insertion involves unsmeared fields and must be
evaluated through some other technique. The two-point contribution follows
from Sec.~\ref{subsec:theory-baryons}.
For an up-quark insertion in the correlation function
\begin{equation}
\begin{split}
   C^{(3)}_\mathrm{conn.}(t_f, t, t_i) =& -\epsilon^{abc}\epsilon^{\bar{a}\bar{b}\bar{c}} S_{\alpha_1\alpha_2 \alpha_3} \bar{S}_{\bar{\alpha}_1 \bar{\alpha}_2 \bar{\alpha}_3}\\
&\quad\quad\Big\langle \big(
({\cal D}_1\Box d)^a_{\alpha_1}
({\cal D}_2\Box u)^b_{\alpha_2}
({\cal D}_3\Box u)^c_{\alpha_3} \big)(t_f)\\&\quad\quad\cdot
\big(\bar{u}\Gamma u\big)(t)
\cdot \big( (\bar d\Box\bar{\cal D}_1)^{\bar{a}}_{\bar\alpha_1}
(\bar u\Box\bar{\cal D}_2)^{\bar{b}}_{\bar\alpha_2}
(\bar u\Box\bar{\cal D}_3)^{\bar{c}}_{\bar\alpha_3}\big)(t_i)
\Big\rangle,
\end{split}
\end{equation}
there are four Wick contractions after quark-field integration.
The up-quark generalized perambulator from Eq.~\ref{eqn:genperam}
is denoted as
\begin{equation}
\widetilde{\tau}_U(t_f,t,t_i)  =
V^\dagger(t_f) M_u^{-1}(t_f,t)\Gamma(t) M_u^{-1}(t,t_i) V(t_i)
\end{equation}
and similarly for the down quark. Since the quarks in the insertion
come in pairs, the connected three-point correlator can be written as the
sum of two two-point correlation functions
with each up-quark perambulator substituted with the corresponding
generalized perambulator in turn;
\begin{equation}
C^{(3)}_u(t_f,t,t_i) = C^{(2)}[\tau_D,\tau_U,\widetilde{\tau}_U](t) +
C^{(2)}[\tau_D,\widetilde{\tau}_U,\tau_U](t).
\end{equation}
The down-quark contribution can be written as
\begin{equation}
C^{(3)}_d(t_f,t,t_i) = C^{(2)}[\widetilde{\tau}_D,\tau_U,\tau_U](t).
\end{equation}
This kind of construction is a generic feature of the three-point functions.

As in the previous two- and three-point correlator examples, the choice
of source and sink operators is independent of the perambulator and
generalized perambulator computations.  These latter generalized
perambulators can be shared among meson three-point correlator
calculations. In addition, for degenerate quarks, only one generalized
perambulator need be computed, independent of the flavor of the quark line
where the current insertion is to occur.
Thus, the overall computation of distinct three-point
correlation functions has a manageable computational cost.
